\documentclass[11pt]{article}
\usepackage[margin=1in]{geometry}
\usepackage{amsmath,amssymb,amsthm}
\usepackage{graphicx}
\usepackage{hyperref}
\usepackage{xcolor}

\newtheorem{theorem}{Theorem}
\newtheorem{lemma}[theorem]{Lemma}
\newcommand{\C}{\mathbb C}
\newcommand{\R}{\mathbb R}

\title{Continuous-Time Stabilization by Compact Memory:\\
A Counterexample Suggested by arXiv:2211.15121}
\author{Automatic mathematics run \texttt{fa\_banach\_001}}
\date{17 August 2026}

\begin{document}
\setlength{\emergencystretch}{3em}
\maketitle

\begin{center}
\fbox{\parbox{0.91\textwidth}{
\textbf{Status.} Candidate counterexample, likely valid.  At the continuity
level of the source's discrete example, a positive essential spectral bound,
a compact nonlinear perturbation, and a quadratically small perturbation do
\emph{not} force continuous-time instability.  The compact map below is not
locally Lipschitz at zero, so stronger regularity variants are not addressed.
}}
\end{center}

\section{Source question}

The source is V. M\"uller, R. Schnaubelt, and Y. Tomilov,
\emph{On growth and instability for semilinear evolution equations: an
abstract approach}, arXiv:2211.15121v3 \cite{source}.  Its Example~5.3
constructs a stable discrete system
\[
 x_{n+1}=(A+K+G)x_n
\]
on an infinite-dimensional Hilbert space, although $r_e(A)>1$, $K$ is a
continuous compact map, and $G(x)=O(\|x\|^2)$ at zero.  Immediately afterward,
on printed page~19, the authors write:
\begin{quote}
``We are short of a similar construction for continuous time. However, we
suspect that an example similar to the above for (1.1) may not exist.''
\end{quote}
The same page recalls the continuous equation and the paper's definition of
stability.  The complete context is reproduced below.

\begin{center}
\includegraphics[width=0.96\textwidth]{figures/open_problem_crop.png}
\end{center}

The precise analogue tested here is whether the zero solution of
\begin{equation}\label{eq:abstract}
 x'(t)=Ax(t)+K(x(t))+G(x(t))
\end{equation}
must be unstable when the essential spectral bound of $A$ is positive, $K$ is
continuous and compact with $K(0)=0$, and $G(x)=O(\|x\|^{1+\delta})$ for some
$\delta>0$.  The answer is no.

\section{Idea of the proof}

The compact perturbation uses only one scalar ``memory'' coordinate.  If $y$
is the component in the infinite-dimensional unstable subspace, the memory
grows at speed $\sqrt{\|y\|}$.  It therefore becomes much larger than $\|y\|$
near zero.  A damping term of size at most $2|a|^2$, hence quadratic in the
full state, can then reverse the growth of $y$ before it leaves a small
neighborhood.

After setting $w=\sqrt{\|y\|}$ and writing $a=u+iv$, the entire stability
question becomes a planar phase estimate.  The memory variable satisfies
$u'=w$, and, with $u$ as the independent variable, the unstable amplitude
satisfies $dw/du=q/2$ for a coefficient $q\in[-1,1]$.  Once
$u^2+v^2\ge w^2$, that region is invariant and $w$ cannot increase.  Soon
afterward $dw/du\le-1/4$, which bounds the total memory accumulation.  This
gives a uniform square-root stability modulus.

\section{Counterexample}

\begin{theorem}\label{thm:main}
Let $H_0$ be any infinite-dimensional complex Hilbert space and put
$H=H_0\oplus\C$ with its Hilbert norm.  There are a bounded generator
$A\in\mathcal B(H)$, a continuous compact map $K:H\to H$, and a continuous
map $G:H\to H$ such that:
\begin{enumerate}
\item $\sigma_e(A)=\{1\}$, so the essential spectral bound of $A$ is $1$;
\item $K(0)=G(0)=0$ and $\|G(x)\|\le2\|x\|^2$ for every $x\in H$;
\item equation~\eqref{eq:abstract} has a unique global classical solution for
every initial state;
\item the zero solution is Lyapunov stable.  More precisely, if
$\eta=\|x(0)\|\le1$, then
\[
 \sup_{t\ge0}\|x(t)\|\le50\sqrt{\eta}.
\]
\end{enumerate}
\end{theorem}

\begin{proof}
Write $x=(y,a)\in H_0\oplus\C$ and $s=\|y\|$.  Define
\begin{align*}
 A(y,a)&=(y,0),\\
 K(y,a)&=(0,\sqrt{s}),\\
 G(y,a)&=\left(-\frac{2|a|^2}{s+|a|^2}\,y,0\right),
\end{align*}
where the last expression is defined to be zero at $(0,0)$.  The operator
$A$ generates the semigroup $(y,a)\mapsto(e^t y,a)$.  Moreover, $A-I$ has an
infinite-dimensional kernel, while $A$ is Fredholm: its kernel and cokernel
are both one-dimensional.  Since $A-\lambda I$ is invertible for
$\lambda\notin\{0,1\}$, it follows directly from the Fredholm definition that
$\sigma_e(A)=\{1\}$.

The map $K$ is continuous.  On a bounded set its image is contained in a
bounded subset of the one-dimensional space $\{0\}\oplus\C$, so it is
precompact; hence $K$ is a compact map.  The apparent singularity in $G$ is
harmless, since
\begin{equation}\label{eq:Gbound}
 \|G(y,a)\|=\frac{2|a|^2s}{s+|a|^2}
 \le2|a|^2\le2\|(y,a)\|^2.
\end{equation}
This estimate also proves continuity at the origin; continuity elsewhere is
immediate (including points $(0,a)$, because the scalar coefficient is
bounded by $2$).

The differential equation is
\begin{equation}\label{eq:system}
 y'=q(s,a)y,\qquad a'=\sqrt{s},\qquad
 q(s,a)=\frac{s-|a|^2}{s+|a|^2},
\end{equation}
with the evident value at $(0,0)$.  Notice that $|q|\le1$.  If $y(0)=0$,
Gronwall's inequality applied to the first equation forces $y\equiv0$, and
then $a$ is constant.  If $y(0)\ne0$, the first equation gives
$y(t)=\exp(\int_0^t q(\tau)\,d\tau)y(0)$, so $y$ never vanishes and its
direction is fixed.  The variables $s=\|y\|$, $u=\operatorname{Re}a$, and
$v=\operatorname{Im}a$ obey a smooth finite-dimensional system on $s>0$:
\[
 s'=qs,\qquad u'=\sqrt{s},\qquad v'=0.
\]
It is locally unique there.  Since $s(0)e^{-t}\le s(t)\le s(0)e^t$ and
$u'=\sqrt{s}$, no component can blow up in finite time.  This proves unique
global well-posedness.

It remains to prove stability.  Assume $s(0)>0$ and set
\[
 w=\sqrt{s},\qquad a=u+iv,\qquad
 q=\frac{w^2-u^2-v^2}{w^2+u^2+v^2}.
\]
Then
\begin{equation}\label{eq:phase}
 w'=\frac q2w,\qquad u'=w,
 \qquad v'=0.
\end{equation}
Thus $u$ is strictly increasing and
\begin{equation}\label{eq:dwdU}
 \frac{dw}{du}=\frac q2\le\frac12.
\end{equation}

First consider any part of the orbit with $u<0$.  Either $u$ stays negative
forever, in which case \eqref{eq:dwdU} already bounds $w$ and all the state
variables, or $u$ first reaches zero with
\[
 w\le w(0)+\frac{|u(0)|}{2}.
\]
Accordingly, start (either at time zero if $u(0)\ge0$, or at this first
zero) with $u_\ast\ge0$ and $w_\ast>0$.  Introduce the convenient initial
bound
\[
 B=w(0)+|u(0)|+|v|;
\]
then $u_\ast\le B$ and $w_\ast\le B$.

Let
\[
 D=u^2+v^2-w^2.
\]
Where $D\ge0$, one has $q\le0$ and
\[
 D'=2uw-qw^2=2uw+\frac{D w^2}{w^2+u^2+v^2}\ge0.
\]
Hence this region is forward invariant.  Before the orbit enters it,
\eqref{eq:dwdU} implies
$w(u)\le w_\ast+(u-u_\ast)/2$.  By
$u=u_\ast+2w_\ast$ we have $u\ge2w_\ast$ and $w\le2w_\ast$, so $D\ge0$.
Thus entry occurs at some $(u_1,w_1)$ satisfying
\begin{equation}\label{eq:entry}
 u_1\le3B,\qquad w_1\le2B.
\end{equation}
After entry, $w$ is non-increasing.  If
$u\ge\sqrt3\,w_1$, then
$u^2+v^2\ge3w^2$ and therefore $q\le-1/2$.  On that part of the orbit,
\[
 \frac{dw}{du}\le-\frac14.
\]
Nonnegativity of $w$ shows that, after $u$ reaches
$\max(u_1,\sqrt3\,w_1)$, it can increase by at most $4w_1$.  Together with
\eqref{eq:entry}, this yields, for all $t\ge0$,
\begin{equation}\label{eq:bounds}
 w(t)\le2B,\qquad |u(t)|\le12B,\qquad |v(t)|\le B.
\end{equation}

Finally let $\eta=\|(y(0),a(0))\|\le1$.  Since
$s(0)\le\eta$, $w(0)\le\sqrt\eta$, and $|a(0)|\le\eta$,
\[
 B\le\sqrt\eta+\sqrt2\,\eta<\frac52\sqrt\eta.
\]
Equations~\eqref{eq:bounds} give
$\|y(t)\|=w(t)^2\le25\eta$ and
$|a(t)|\le13B<\frac{65}{2}\sqrt\eta$.  Consequently
\[
 \|(y(t),a(t))\|<50\sqrt\eta
\]
for all $t\ge0$.  The case $s(0)=0$ was already seen to be stationary.  Given
$\varepsilon>0$, choosing
$\eta<\min\{1,(\varepsilon/50)^2\}$ proves Lyapunov stability of zero.
\end{proof}

\section{Verification and limitations}

The proof is self-contained.  The separate verifier report audits compactness,
essential spectrum, well-posedness, and every phase-plane inequality.  A small
deterministic Runge--Kutta script in \texttt{code/simulate\_reduced.py}
checks representative initial phases and scales; it is a sanity check, not
part of the proof.

The construction depends essentially on the sublinear memory law
$\sqrt{\|y\|}$.  The compact map $K$ is continuous (indeed, $1/2$-H\"older in
the radial variable) but is not locally Lipschitz or Fr\'echet differentiable
at zero.  Therefore this packet does not answer a version imposing local
Lipschitz, $C^1$, or specified differentiability hypotheses on the compact
perturbation.  It also does not contradict the source's positive instability
theorems, which concern compact nonlinearities without the additional
noncompact metrically small term $G$ or invoke stronger local hypotheses.

\section{Novelty check and review recommendation}

A bounded search through 17 August 2026 covered the run's four lightweight
indexes, the source text and bibliography, the exact quoted sentence, the
paper title and author names, and web/arXiv-oriented combinations of
``continuous time,'' ``compact perturbation,'' ``metrically small,''
``essential spectral bound,'' and ``stable zero solution.''  It found the
source paper and general linearized-instability literature, but no later
paper giving this construction or explicitly resolving the quoted suspicion.
The novelty claim is therefore limited to this explicit continuity-level
counterexample.

Human review should focus on two points: whether ``similar construction'' is
intended to include compact maps that are merely continuous at zero, as in the
discrete setup; and the uniqueness argument at the non-Lipschitz equilibrium.
The latter is supplied directly by the bounded scalar coefficient in the
$y$ equation, not by a general Banach-space existence theorem.

\begin{thebibliography}{9}
\bibitem{source}
V. M\"uller, R. Schnaubelt, and Y. Tomilov,
\emph{On growth and instability for semilinear evolution equations: an
abstract approach}, Mathematische Annalen 389 (2024), 3885--3933;
arXiv:2211.15121v3.
\end{thebibliography}

\end{document}
